\section{Exactification in dimension three}
\label{sec:exactification}

This section converts the approximate eigenfunctions of
Proposition~\ref{prop:wave-superposition} into exact eigenfunctions of nearby
ordinary Riemannian metrics.  The conversion preserves all pointwise signs
and leaves a prescribed outer region unchanged.  We first correct the
equation using a scalar conductivity and a scalar density.  A positive
change of the dependent variable then enforces the relation between these
two quantities that a three-dimensional metric requires.  A correction in
a clean annulus permits this change of dependent variable to be cut off.

Throughout this section, \(g\) is fixed while the integer \(n\) tends to
infinity, and \(\Lambda=\Lambda_n=n(n+2)\).  Norms are taken with respect to
fixed smooth reference norms; replacing them by the norms of \(g\) changes
only constants.  Covariant derivatives and pointwise contractions in the
formulas below use \(g\), unless another metric is displayed.  All constants
may depend on this fixed metric and on the finite derivative orders under
consideration, but not on \(n\).

\subsection{A scalar correction with weighted estimates}

\begin{lemma}[Scalar correction]
\label{lem:exact-quotients}
Let \((M,g)\) be a fixed closed smooth Riemannian manifold.  Fix an integer
\(B\ge1\), and let \(h\ge0\) and \(D_*>0\) be integers.  Suppose that, for
all sufficiently large integers \(n\), a real smooth function \(U=U_n\)
and a positive continuous function \(W=W_n\) satisfy
\begin{align}
 |U|+n^{-1}|\nabla U|&\ge n^{-B}W,
 \label{eq:exact-jet-hypothesis}\\
 |\nabla^jU|&\le C_j n^{j+4}W,
 &|\nabla^j\mathcal R|&\le C_jn^{4-D_*}W
 \quad (0\le j\le h+2),
 \label{eq:exact-weighted-hypothesis}
\end{align}
where \(\mathcal R=(\Delta_g+\Lambda)U\) and the constants \(C_j\) are
independent of \(n\).  Define
\begin{equation}
 \begin{aligned}
 \mathcal D&=n^2U^2+|\nabla U|^2,\qquad
 b=1-\frac{\mathcal R U}{\mathcal D},\\
 s&=1-\Lambda^{-1}\left[
 \frac{\mathcal R n^2U}{\mathcal D}
 -\Div_g\left(\frac{\mathcal R}{\mathcal D}\nabla U\right)\right].
 \end{aligned}
 \label{eq:exact-scalar-correction}
\end{equation}
Then \(b,s\) are smooth, and, with
\(\Delta_b f=\Div_g(b\nabla f)\),
\begin{equation}
 \Delta_bU+\Lambda sU=0.
 \label{eq:exact-weighted-equation}
\end{equation}
They equal one on every open set where \(\mathcal R=0\).  Moreover, if
\[
 E=2B+20,\qquad Q_h=(E+1)(h+2)+10,
\]
then
\begin{equation}
 \|b-1\|_{C^h}+\|s-1\|_{C^h}
 \le C_h n^{-D_*+Q_h}.
 \label{eq:exact-quotient-bound}
\end{equation}
In particular, for any prescribed integer \(P_1\ge1\), choosing
\(D_*\ge P_1+Q_h+2\) makes the left side at most \(n^{-P_1}\) for all
sufficiently large \(n\).  In this case \(b\) and \(s\) are positive for
all sufficiently large \(n\).
\end{lemma}

\begin{proof}
The lower bound in \eqref{eq:exact-jet-hypothesis} gives
\begin{equation}
 \mathcal D\ge \tfrac12 n^{2-2B}W^2>0.
 \label{eq:exact-denominator-lower}
\end{equation}
Thus the formulas in \eqref{eq:exact-scalar-correction} define smooth
functions, regardless of the regularity of \(W\).  Put
\(f=\mathcal R/\mathcal D\).  The product rule gives
\begin{align*}
 \Delta_bU+\Lambda sU
 &=\Delta_gU-U\Div_g(f\nabla U)-f|\nabla U|^2
   +\Lambda U-n^2fU^2+U\Div_g(f\nabla U)\\
 &=\mathcal R-f\mathcal D=0.
\end{align*}
The assertion on the set where the residual vanishes follows from the
same formulas.

We give a quantitative derivative estimate to make clear why an
exponentially small weight causes no loss.  For \(0\le j\le h+1\),
the product rule and \eqref{eq:exact-weighted-hypothesis} yield
\[
 |\nabla^j\mathcal D|\le C_j n^{j+10}W^2.
\]
For \(j\ge1\), each term in a derivative of
\(\mathcal D^{-1}\) is a constant multiple of
\[
 \mathcal D^{-p-1}
 (\nabla^{a_1}\mathcal D)\cdots
 (\nabla^{a_p}\mathcal D),
 \qquad a_i\ge1,\quad a_1+\cdots+a_p=j,
\]
with the appropriate tensor contractions.  By
\eqref{eq:exact-denominator-lower}, its norm is bounded by
\[
 C_jn^{j+10p+(2B-2)(p+1)}W^{-2}
 =C_jn^{j+(2B+8)p+2B-2}W^{-2}.
\]
The same conclusion for \(j=0\) follows directly from
\eqref{eq:exact-denominator-lower}.  Since \(p\le j\), we may use the
convenient weaker bound
\begin{equation}
 |\nabla^j(\mathcal D^{-1})|
 \le C_jn^{E(j+1)}W^{-2}
 \qquad (0\le j\le h+1).
 \label{eq:exact-inverse-derivatives}
\end{equation}

A derivative of \(b-1\) consists of a derivative of \(\mathcal R\),
a derivative of \(U\), and a derivative of \(\mathcal D^{-1}\).
The first two factors contribute one power of \(W\) each and the last
contributes \(W^{-2}\).  For \(j\le h\), the resulting bound is at most
\[
 C_jn^{-D_*+8+j+E(j+1)}.
\]
The term \(\Lambda^{-1}\mathcal Rn^2U/\mathcal D\) has the same
bound, because \(n^2/\Lambda\le1\).  Finally, a derivative of order
\(j\) of
\(\Lambda^{-1}\Div_g((\mathcal R/\mathcal D)\nabla U)\)
uses derivatives of \(\mathcal R\) through order \(j+1\) and of
\(U\) through order \(j+2\).  Applying
\eqref{eq:exact-inverse-derivatives}, and using
\(\Lambda^{-1}\le n^{-2}\), bounds it by
\[
 C_jn^{-D_*+E(j+2)+j+8}.
\]
Both displayed exponents are bounded above by \(-D_*+Q_h\).
This proves \eqref{eq:exact-quotient-bound}.  At no point have we
differentiated \(W\); its powers cancel in every differentiated
quotient.  The final conclusions follow by taking \(n\) large enough
to absorb the fixed constant and to make the \(C^0\) error smaller
than \(1/2\).
\end{proof}

\subsection{A positive change of the dependent variable}

The corrected coefficients $b,s$ may vary independently, whereas the
conductivity and density of an ordinary metric satisfy a determinant
relation.  In dimension three the relation takes the following form.

\begin{lemma}[Metric conductivity]
\label{lem:metric-conductivity}
Let \((M,g)\) be a smooth three-dimensional Riemannian manifold and
let \(A\) be a smooth positive \(g\)-self-adjoint endomorphism of
\(TM\).  Put \(d=\det A\), with determinant taken in any
\(g\)-orthonormal frame, and define
\begin{equation}
 g_A(X,Y)=g(dA^{-1}X,Y).
 \label{eq:exact-metric-realization}
\end{equation}
Then \(g_A\) is a smooth Riemannian metric and
\begin{equation}
 \dd V_{g_A}=d\,\dd V_g,\qquad
 \Delta_{g_A}f=d^{-1}\Div_g(A\nabla f).
 \label{eq:exact-metric-operator}
\end{equation}
\end{lemma}

\begin{proof}
In a \(g\)-orthonormal frame the metric endomorphism is
\(C=dA^{-1}\).  It is positive and symmetric, and in dimension three
\[
 \det C=\frac{d^3}{\det A}=d^2.
\]
Thus its positive relative volume density is \(d\), and its inverse
endomorphism is \(C^{-1}=A/d\).  The energy identity is consequently
\[
 \int_M g_A(\nabla_{g_A}f,\nabla_{g_A}h)\,\dd V_{g_A}
 =\int_M g(A\nabla f,\nabla h)\,\dd V_g
\]
for compactly supported smooth \(f,h\).  Integration by parts proves
the operator identity in \eqref{eq:exact-metric-operator}.
\end{proof}

To see which equation the positive factor must satisfy, suppose that
$b,s>0$ and $U$ satisfy \eqref{eq:exact-weighted-equation}.  For any
positive smooth function $v$, define
\[
 q_v=bv^2,\qquad
 d_v=sv^2+\Lambda^{-1}v\Delta_bv.
\]
The product rule gives
\begin{align}
 q_v\nabla(U/v)&=b(v\nabla U-U\nabla v),\notag\\
 \Div_g\bigl(q_v\nabla(U/v)\bigr)+\Lambda d_v(U/v)&=0.
 \label{eq:exact-division-identity}
\end{align}
Indeed, the divergence of the displayed flux is
$v\Delta_bU-U\Delta_bv=-\Lambda d_v(U/v)$.
To realize this equation by the scalar conductivity $q_v\Id$, we impose
$d_v=q_v^3$.  After division by $v>0$, this condition is
\[
 -\Delta_bv+\Lambda b^3v^5-\Lambda sv=0.
\]
The next lemma solves this equation with $v$ close to one.  The solution
need not equal one in the prescribed exterior; the annular correction
below will allow us to cut it off there.

\begin{lemma}[A positive conjugating factor]
\label{lem:exact-positive-factor}
Let \((M,g)\) be a fixed closed smooth three-dimensional Riemannian
manifold.  Given integers \(k_2\ge0\) and \(P_2\ge1\), choose an
even integer \(l\ge4\) with \(l>k_2+3/2\), and set
\[
 k_1=l,\qquad P_1=P_2+8.
\]
Suppose that smooth functions \(b=b_n\) and \(s=s_n\) satisfy
\begin{equation}
 \|b-1\|_{C^{k_1}}+\|s-1\|_{C^{k_1}}
 \le n^{-P_1}.
 \label{eq:exact-coefficient-smallness}
\end{equation}
For all sufficiently large \(n\), there is a positive smooth function
\(v=v_n\) solving
\begin{equation}
 -\Delta_bv+\Lambda b^3v^5-\Lambda sv=0
 \label{eq:exact-semilinear}
\end{equation}
and satisfying
\begin{equation}
 \|v-1\|_{C^{k_2}}\le n^{-P_2}.
 \label{eq:exact-factor-bound}
\end{equation}
The required threshold for \(n\) may depend on the fixed manifold and
the displayed finite orders.
\end{lemma}

\begin{proof}
Let \(H^a=H^a(M,g)\) denote the usual \(L^2\) Sobolev spaces.  On the
fixed manifold, the inverse
\[
 L_\Lambda^{-1}=(-\Delta_g+4\Lambda)^{-1}:H^{l-2}\longrightarrow H^l
\]
has norm bounded independently of \(\Lambda\ge1\).  Indeed, if
\(\mu\ge0\) is an eigenvalue of \(-\Delta_g\), the multiplier for
this map in the spectral Sobolev norms is
\[
 \frac{1+\mu}{\mu+4\Lambda}\le1.
\]
The spectral and chart Sobolev norms are equivalent with constants
depending on \(g\) and \(l\), not on \(\Lambda\).  This equivalence
follows from the compact-manifold elliptic estimate and the spectral
theorem; see \cite[Theorems 15.7 and 16.1]{DyatlovNotes}.

Write \(v=1+h\).  Equation~\eqref{eq:exact-semilinear} becomes
\begin{equation}
 L_\Lambda h
 =\Div_g((b-1)\nabla h)
  -\Lambda\bigl(b^3(1+h)^5-s(1+h)-4h\bigr).
 \label{eq:exact-fixed-point-equation}
\end{equation}
Let \(T_n(h)\) be the result of applying \(L_\Lambda^{-1}\) to the
right side.  The polynomial in parentheses has constant term
\(b^3-s\) and linear coefficient
\[
 5b^3-s-4=5(b^3-1)-(s-1).
\]
Its remaining terms have degree at least two in \(h\).

Sobolev multiplication and \eqref{eq:exact-coefficient-smallness}
therefore imply, on the closed \(H^l\) ball of radius \(\delta\le1\),
\begin{align*}
 \|T_n(0)\|_{H^l}&\le C\Lambda n^{-P_1},\\
 \|T_n(h)-T_n(\widetilde h)\|_{H^l}
 &\le C\Lambda(n^{-P_1}+\delta)
       \|h-\widetilde h\|_{H^l}.
\end{align*}
Here the divergence term is bounded from \(H^l\) to \(H^{l-2}\)
by \(C\|b-1\|_{C^l}\); the harmless factor \(\Lambda\ge1\)
allows it to be included in the displayed estimate.  The multiplication
and embedding statements used here hold in dimension three because
\(l\ge4\) and \(l>k_2+3/2\); see
\cite[Section~1, Exercise~6, and Proposition~3.3]{TaylorSobolev}.

Set \(Q=P_2+4\) and \(\delta=n^{-Q}\).  Since \(P_1=Q+4\) and
\(\Lambda\le3n^2\) for \(n\ge1\),
\[
 \|T_n(0)\|_{H^l}\le Cn^{-Q-2},\qquad
 \operatorname{Lip}(T_n)\le Cn^2(n^{-P_1}+n^{-Q})=o(1).
\]
For all sufficiently large \(n\), the map preserves this closed ball
and is a contraction there.  Its fixed point satisfies
\(\|h\|_{H^l}\le n^{-Q}\).  Sobolev embedding then gives
\(\|h\|_{C^{k_2}}\le Cn^{-Q}\le n^{-P_2}\), and also
\(\|h\|_{C^0}<1/2\), so \(v=1+h\) is positive.

The coefficients \(b,s\) are smooth and \(b\) is uniformly positive
for large \(n\).  The polynomial right side of
\(\Delta_bv=\Lambda b^3v^5-\Lambda sv\) lies initially in \(H^l\).
Elliptic regularity raises the regularity of \(v\) to \(H^{l+2}\);
iteration gives \(v\in C^\infty\).  Only smoothness, not a bound
uniform in \(n\) on every derivative at once, is needed in this
bootstrap; see \cite[Theorems 15.1 and 15.7]{DyatlovNotes}.
\end{proof}

\subsection{Localization and annular determinant repair}

\begin{proposition}[Exactification with an unchanged exterior]
\label{prop:exactification}
Let \(g\) be a fixed smooth metric on \(\Sph^3\), and fix
\(0<t_2<t_3<1\).  Assume that \(g=g_*\) on \(\{r\ge t_2\}\), where
\(g_*\) is the unit round metric and
\(r=(x_1^2+x_2^2)^{1/2}\).  Fix an integer \(B\ge1\).

For every pair of integers \(k\ge0\), \(P\ge1\), there exist finite
integers \(m_*,D_*\), depending only on \(k,P,B\), with the following
property.  Suppose that, for all sufficiently large integers \(n\),
there are real smooth functions \(U_n\) and positive continuous
functions \(W_n\) such that
\begin{gather}
 |U_n|+n^{-1}|\nabla U_n|\ge n^{-B}W_n,
 \label{eq:exact-proposition-lower}\\
 |\nabla^jU_n|\le C_jn^{j+4}W_n,\qquad
 |\nabla^j((\Delta_g+\Lambda_n)U_n)|
 \le C_jn^{4-D_*}W_n
 \quad(0\le j\le m_*),
 \label{eq:exact-proposition-upper}\\
 U_n=w_n=\Re(x_1+ix_2)^n\quad\hbox{on }\{r\ge t_2\},
 \label{eq:exact-proposition-exterior}
\end{gather}
where all \(C_j\) are independent of \(n\).  Then, for all sufficiently
large \(n\), there are a smooth Riemannian metric \(\widehat g_n\)
and a positive smooth function \(\widetilde v_n\) on \(\Sph^3\) such
that, with \(u_n=U_n/\widetilde v_n\),
\begin{gather}
 (\Delta_{\widehat g_n}+\Lambda_n)u_n=0,
 \qquad \|\widehat g_n-g\|_{C^k}\le n^{-P},
 \label{eq:exact-final-equation}\\
 \widehat g_n=g_*,\quad \widetilde v_n=1,\quad u_n=w_n
 \quad\hbox{on }\{r\ge t_3\}.
 \label{eq:exact-final-exterior}
\end{gather}
In particular \(u_n\) is nonzero and has exactly the same pointwise
signs and zero set as \(U_n\).  The threshold for \(n\) may depend
on the fixed input data and on the constants in
\eqref{eq:exact-proposition-upper}.
\end{proposition}

\begin{proof}
We specify the finite orders before taking \(n\) large.  Choose
\begin{equation}
 \begin{gathered}
 k_2=k+2,\qquad P_2=P+k+10,\\
 l\ge4\text{ even},\quad l>k_2+3/2,\qquad
 k_1=l,\quad P_1=P_2+8,\\
 m_*=k_1+2,\qquad
 D_*\ge P_1+(2B+21)(k_1+2)+12.
 \end{gathered}
 \label{eq:exact-parameter-budget}
\end{equation}
These choices depend only on \(k,P,B\).  We suppress the subscript
\(n\) in the construction.

Apply Lemma~\ref{lem:exact-quotients} with \(h=k_1\) to obtain
positive smooth coefficients \(b,s\) satisfying
\eqref{eq:exact-weighted-equation} and
\eqref{eq:exact-coefficient-smallness}.  Since \(g\) is round and
\(U=w_n\) on \(\{r\ge t_2\}\), their residual vanishes there,
and \(b=s=1\) there.  Apply
Lemma~\ref{lem:exact-positive-factor} to obtain a positive smooth
solution \(v\) of \eqref{eq:exact-semilinear} satisfying
\eqref{eq:exact-factor-bound}.

Choose numbers \(t_2<a<b_0<t_3\), and fix a smooth cutoff \(\chi\)
with \(0\le\chi\le1\), which is one on \(\{r\le a\}\) and zero on
\(\{r\ge b_0\}\).
Such a cutoff is smooth on the sphere because it is constant near
\(r=0\).  Set
\begin{equation}
 \begin{aligned}
 \widetilde v&=1+\chi(v-1),\qquad u=U/\widetilde v,\\
 q&=b\widetilde v^2,\qquad
 d=s\widetilde v^2+\Lambda^{-1}\widetilde v\Delta_b\widetilde v.
 \end{aligned}
 \label{eq:exact-conjugation}
\end{equation}
For all sufficiently large \(n\), \(\widetilde v\) is positive and
\(\|\widetilde v-1\|_{C^{k+2}}\le Cn^{-P_2}\).  Applying the division
identity \eqref{eq:exact-division-identity} with
\(v=\widetilde v\) gives
\begin{equation}
 \Div_g(q\nabla u)+\Lambda d u=0.
 \label{eq:exact-conjugated-equation}
\end{equation}

The coefficient estimates and the two extra derivatives included in
\(k_2=k+2\) give
\begin{equation}
 \|q-1\|_{C^k}+\|d-1\|_{C^k}\le Cn^{-P_2}.
 \label{eq:exact-density-estimates}
\end{equation}
For the term \(\Delta_b\widetilde v\), one uses
\(\Delta_b\widetilde v=b\Delta_g\widetilde v+
g(\nabla b,\nabla\widetilde v)\), so the required derivatives of
\(b\) are available because \(k_1\ge k+1\).
In particular \(q,d>0\) for large \(n\).

On the open set \(\{r<a\}\), the cutoff is identically one, and
\eqref{eq:exact-semilinear} gives
\[
 d=sv^2+\Lambda^{-1}v\Delta_bv=b^3v^6=q^3.
\]
Where \(r\ge b_0\), we have \(\widetilde v=1\) and \(b=s=1\),
so \(q=d=1\).  Consequently the defect \(d-q^3\) is supported in
the compact annulus \(\{a\le r\le b_0\}\), strictly inside
\(\{t_2<r<t_3\}\).  Without this defect, the choice \(A=q\Id\)
would already have the required determinant \(d\).  We correct that
determinant on the two-plane perpendicular to \(\nabla u\).

We first justify that this plane and its derivatives are controlled
throughout the transition annulus.  The round gradient of \(w_n\)
satisfies
\begin{equation}
 |\nabla w_n|_{g_*}^2
 =n^2(r^{2n-2}-w_n^2)
 \ge n^2r^{2n-2}(1-r^2).
 \label{eq:exact-annular-gradient}
\end{equation}
To verify the equality, the ambient Euclidean gradient has squared
norm \(n^2r^{2n-2}\), while its normal component on the unit sphere
is \(nw_n\), by homogeneity.  Since \(U=w_n\) here and
\(\widetilde v\) is close to one in \(C^1\),
\[
 |\nabla u|
 \ge \widetilde v^{-1}nr^{n-1}\sqrt{1-r^2}
       -\widetilde v^{-2}r^n|\nabla\widetilde v|
 \ge c n r^n>0
 \quad(t_2<r<t_3)
\]
for large \(n\).

For the derivative estimates, write locally
\(x_1+ix_2=re^{i\theta}\) and set
\[
 \alpha_n=\frac{\cos(n\theta)}r\nabla r
             -\sin(n\theta)\nabla\theta.
\]
The displayed expression is globally defined on the annulus: the
trigonometric functions are single-valued and
\(\dd\theta=(-x_2\dd x_1+x_1\dd x_2)/r^2\) is a smooth one-form
there.  On the fixed annulus,
\[
 |\alpha_n|^2=r^{-2}-\cos^2(n\theta)\ge t_3^{-2}-1>0,
 \qquad |\nabla^j\alpha_n|\le C_j n^j.
\]
Moreover
\[
 \nabla u=nr^n t_n,\qquad
 t_n=\widetilde v^{-1}\alpha_n
      -n^{-1}\cos(n\theta)\widetilde v^{-2}\nabla\widetilde v.
\]
It follows that \(|t_n|\ge c>0\) and
\(|\nabla^jt_n|\le C_jn^j\) for \(0\le j\le k\).  The
orthogonal projection
\[
 P_\perp=\Id-\frac{t_n\otimes t_n^\flat}{|t_n|^2}
\]
is therefore smooth and satisfies
\begin{equation}
 |\nabla^jP_\perp|\le C_jn^j\qquad(0\le j\le k).
 \label{eq:exact-projection-estimates}
\end{equation}
Here \(t_n^\flat=g(t_n,\cdot)\), so
\((t_n\otimes t_n^\flat)X=g(t_n,X)t_n\).
To obtain \eqref{eq:exact-projection-estimates}, differentiate the
quotient defining \(P_\perp\); its denominator is uniformly bounded
away from zero, and the total derivative order of the numerator factors
in every term is at most \(j\).  This estimate shows explicitly that
only \(k\) powers of \(n\) are lost in the required projection norms.

Define a positive symmetric endomorphism by
\begin{equation}
 A=q\left[\Id+\left(\sqrt{d/q^3}-1\right)P_\perp\right]
 \quad\hbox{on }\{t_2<r<t_3\},\qquad
 A=q\Id\quad\hbox{elsewhere}.
 \label{eq:exact-annular-conductivity}
\end{equation}
The coefficient of \(P_\perp\) vanishes on neighborhoods of both
boundary components of this annulus, because \(d=q^3\) there.
Thus the two formulas agree on open collars and define a globally
smooth endomorphism.  Its eigenvalue along \(\nabla u\) is \(q\),
and its two transverse eigenvalues are both
\(q\sqrt{d/q^3}\).  Hence, throughout the annulus,
\[
 A\nabla u=q\nabla u,\qquad \det A=d.
\]
These identities also hold outside it, where \(A=q\Id\) and
\(d=q^3\).

By \eqref{eq:exact-density-estimates}, smooth composition gives
\(\|\sqrt{d/q^3}-1\|_{C^k}\le Cn^{-P_2}\).
The product rule and \eqref{eq:exact-projection-estimates} consequently
give
\[
 \|A-\Id\|_{C^k}\le Cn^{k-P_2}=Cn^{-P-10}.
\]
Define \(\widehat g=g_A\) by
\eqref{eq:exact-metric-realization}.  The map
\((d,A)\mapsto dA^{-1}\) is smooth near \((1,\Id)\), so
\eqref{eq:exact-density-estimates} and this last estimate imply
\[
 \|\widehat g-g\|_{C^k}\le Cn^{-P-10}\le n^{-P}
\]
for all sufficiently large \(n\).  Combining
Lemma~\ref{lem:metric-conductivity} with
\eqref{eq:exact-conjugated-equation} proves
\(\Delta_{\widehat g}u=-\Lambda u\).  On \(\{r\ge t_3\}\),
\(A=\Id\), \(\widetilde v=1\), and \(U=w_n\), which proves
\eqref{eq:exact-final-exterior}.  The lower bound
\eqref{eq:exact-proposition-lower} excludes \(U\equiv0\), and
division by the positive function \(\widetilde v\) preserves all
pointwise signs and zeros.
\end{proof}

\begin{remark}
The wave weight \(W_n=e^{n\phi}+r^n\) satisfies the
positivity and continuity hypotheses of Proposition~\ref{prop:exactification}.
The orders in \eqref{eq:exact-parameter-budget} are chosen before the
finite phase and amplitude orders needed for the wave construction,
and all these choices precede the final choice of \(n\).  Increasing
the threshold for \(n\) then absorbs every constant depending on the
fixed profile or those finite orders.  Consequently the conclusion
gives arbitrary closeness in any prescribed finite collection of
smooth norms.

Hence the sign bound \eqref{eq:wave-sign-lower} survives with the same
constant, while the unchanged exterior in
\eqref{eq:exact-final-exterior} leaves room for the spectral
perturbation of the next section.
\end{remark}
